\subNsection{1}{Perspektiven}{Kristian Popov} \label{6_3_Perspektiven} \label{6_2_Perspektiven}

Die modulare Architektur erlaubt eine flexible Erweiterung der Anwendung. Künftige Perspektiven gliedern sich in die folgenden Schwerpunkte:

\begin{itemize}
    \item \textbf{Weitere Lernbereiche und Topics:} Neue Themen, etwa in Form von Fachmodulen, die von anderen Professoren beigesteuert werden, lassen sich in Zukunft über den zentralen Modulkatalog dynamisch in Navigation und Routing integrieren.
    \item \textbf{Aufgabenverwaltung und -varianten:} Lehrende können Aufgaben und Vorlagen direkt über die Oberfläche anlegen, automatisch Varianten erzeugen sowie Quizfragen und Lernkarten ergänzen.
    \item \textbf{Mobile Nutzung:} Die Benutzeroberfläche, insbesondere die Übergangstabellen und Zustandsdiagramme, für kleine Bildschirme und Tablets optimieren.
    \item \textbf{Java-Zeilensperrung:} Abgrenzung bearbeitbarer Code-Bereiche im Editor durch Marker, um nicht zu bearbeitende Zeilen schreibgeschützt zu halten.
    \item \textbf{Gamification und LMS-Anbindung:} Motivationselemente wie Streaks oder Abzeichen sowie eine Moodle-Anbindung zur zentralen Lernstandsauswertung.
\end{itemize}

Neue Lernbereiche, Aufgabenvarianten und Editor-Sperren würden sich vollständig in die bestehende Client-Architektur integrieren lassen. Zentrale Auswertungen, Nutzerkonten oder eine direkte Moodle-Anbindung hingegen würden ein Backend erfordern.
